\section{Introduction}
% \wyc{The outline of the paragraph before proposing \tool:
% 1. Constraint solving is fundamental for addressing problems expressed in various logics.

% 2. Briefly explain existing works on solving constraints. 

% 3. Introduce the important direction (\ie, select important variables and promising values) that existing methods are not aware of.
% }
\SMT{} solving is a fundamental technology in modern computer science, providing a core reasoning engine for diverse software engineering tasks, including software testing~\cite{cadar2008klee}, program verification~\cite{ball2004slam}, and synthesis~\cite{reynolds2019refutation}. Despite the remarkable power of modern solvers like Z3~\cite{demoura2008z3} and CVC5~\cite{cvc5}, a fundamental and persistent challenge remains: the scalability bottleneck of SMT solving. This bottleneck becomes particularly acute when dealing with complex, real-world constraints involving non-linear arithmetic, high-dimensional variables, and mixed logics, frequently leading to timeouts and limiting the effectiveness of client applications~\cite{cadar2013symbolic}.
% \lz{Compared with the original problem, abstract the problem to a more advanced level. The three types of work represent different methods to solve this problem, and our method is a new improvement made in one of these categories.}

To address this challenge, existing methods largely focus on two primary strategies.
% \lz{Reorganized to focus on the two most relevant categories for SMT solving scalability, removing path-level optimization as suggested by wyc reviewer.} 
The first direction is \textit{solver-level enhancement}, which focuses on improving the internal mechanisms of the solver itself. This ranges from meta-level strategies like adaptive parameter tuning (\eg{}, SymTuner~\cite{chaSymTunerMaximizingPower2022}) to embedding learning directly into the solver's core via neural branching heuristics~\cite{kurin2020improving}. 
The second strategy is constraint-level simplification, which operates on the formula itself before it is passed to the solver. This strategy encompasses a range of techniques that largely rely on fixed policies, from general-purpose syntactic methods like rewrite-based normalization~\cite{willsey2021egg}, to sophisticated, specialized algorithms designed for specific theories like array constraints~\cite{liu2022optimal, shuai2021type}.
% \lz{focus on 'fixed-strategy' methods, contrast with COMPASS's 'learned-policy' approach, sharpening our core novelty argument.}

While each of these strategies has proven valuable, they possess fundamental limitations regarding the guidance gap---the lack of high-level, semantic insight to guide reasoning.
% \lz{Updated to focus on the two remaining strategies after removing path-level optimization, maintaining logical flow to establish research gap.} 
Solver-level enhancement, while powerful, does not reduce the intrinsic semantic complexity of the constraints themselves and is often tightly coupled with solver internals. This leaves constraint-level simplification as the most direct approach to tackling a formula's complexity. However, the techniques in this category, whether syntactic or specialized, are almost universally driven by fixed strategies and brittle, hand-crafted heuristics. They fundamentally lack a learned, adaptive policy to understand a formula's semantic structure and make strategic, variable-level simplification decisions.
This analysis reveals a clear and pressing research gap: a systematic, solver-agnostic method to learn where to simplify (\ie, identifying high-impact variables) and how to simplify (\ie, proposing semantically promising values).

To bridge this gap, we present \tool (\textbf{CO}llaborative \textbf{M}ethod with \textbf{P}olicy-learning \textbf{A}nd \textbf{S}emantic \textbf{S}ynthesis)\footnote{Analogous to a compass—where the needle (\RL{}) sets the strategic direction and the dial (\LLM{}) provides tactical detail—\tool integrates these roles to guide \SMT{} solvers efficiently through vast and complex search spaces.}, a novel framework that reframes constraint simplification as a sequential decision-making problem. The core idea of \tool is to use a dual-agent AI architecture that performs strategic, variable-level concretization by leveraging the distinct strengths of each component. This architecture decomposes the complex simplification task: a strategic \RL{} agent, whose policy is continuously adapted, decides what to simplify (variable selection), while a pre-trained \LLM{}, serving as a stable source of general semantic knowledge, determines how to simplify (value suggestion).
% \lz{Here, we consider summarizing some explanations directly from the method section to illustrate why only the policy of RL is updated. However, given that we have previously mentioned that the RL agent  (whose policy is continuously adapted) and that the LLM (serves as a stable source of general semantic knowledge), we will not provide excessive additional elaboration here.} 
The \RL{} agent's policy is subsequently refined based on feedback gathered from this interactive process, which is quantified and guided by our hybrid reward system. This system provides a rich learning signal by combining sparse, ground-truth feedback from the \SMT{} solver, dense heuristic guidance from pre-trained predictive models, and history-based penalties that discourage re-exploring failed paths. Through this feedback loop, \tool learns to recognize a formula's semantic structure and applies policy-guided simplifications that significantly improve downstream solving efficiency.

To validate our approach, we evaluated \tool on \QFLIA{} and \QFNIA{} theories from SMTimer and \SMTCOMP{} benchmarks, using four state-of-the-art solvers: Z3~\cite{demoura2008z3}, CVC5~\cite{cvc5}, MathSAT~\cite{mathsat5}, and BVParti~\cite{bvparti2025}. \tool demonstrates consistent and substantial improvements. Specifically, on SMTimer with Z3, \tool increased the number of solved instances from 72 to 94 (+30.6\%) and reduced the average solving time from 561s to 214s (2.6$\times$ speedup). The improvements extend across all tested solvers: CVC5 achieved a +57.9\% increase in solved instances, MathSAT +244.8\%, and BVParti +850\%. 
Moreover, on the SMT-COMP benchmarks, \tool achieved a +4.6\% improvement on \QFLIA{}, and on \QFNIA{}, it enabled Z3 to solve 324 instances that previously timed out within the same 1200s budget. \wyc{I think the result introduction should be consistent, \eg, \tool achieved a +4.6\% (\resp, +xxx\%) improvement on \QFLIA{} (\resp, \QFNIA{}).}
These results demonstrate that policy-driven, semantics-aware concretization can turn otherwise intractable formulas into solvable ones while preserving solver independence.

The main contributions of this paper are as follows:
\begin{itemize}
  \item We formalize \SMT{} constraint simplification as a sequential decision-making problem and present \tool, a novel AI-assisted framework that employs a dual-agent architecture that decouples strategic variable selection (\RL{}) from tactical value suggestion (\LLM{}).

  \item We apply \tool to four existing SOTA solvers without modifying their internals. A complexity-guided, time-bucket-aware routing mechanism selectively invokes the AI-assisted simplification under tight budgets, enabling deployment-oriented cost–benefit trade-offs suitable for production environments.

  \item We conduct comprehensive evaluations demonstrating \tool's effectiveness across challenging constraint sets and solver architectures. \tool achieves substantial improvements: 30.6\% more solved instances and 2.6$\times$ speedup with Z3, with consistent gains across SOTA solvers. Ablation studies confirm that the \RL{}-\LLM{} synergy substantially outperforms individual components and baselines.
\end{itemize}

The remainder of this paper is organized as follows: Section~\ref{sec:background} provides the necessary background. Section~\ref{sec:motivation-and-problem} presents the motivating example and formalizes the task as an \MDP{}. Section~\ref{sec:methodology} details \tool. Section~\ref{sec:evaluation} reports our extensive experimental evaluation. Section~\ref{sec:discussion} discusses limitations, and Section~\ref{sec:related} reviews related work, before we conclude in Section~\ref{sec:conclusion}. \wyc{Can be removed when exceeding page limit.}
